\chapter{Object Storage (Feature Branch)}
\label{chap:object-storage}

\begin{quote}
\textbf{Scope.} Everything in this chapter describes the
\verb|feat/object-storage| branch. It is \emph{not} merged into \verb|main|; the
\verb|main| database connector documented in Chapter~\ref{chap:database} is
filesystem-only. This chapter is included so the reference is complete, and each
feature-branch-only element is called out as such.
\end{quote}

\section{Motivation}

LanceDB can persist to cloud object storage (S3, GCS, Azure Blob) as easily as to
a local directory --- the on-disk Lance format is the same, only the
\verb|object_store| backend differs. Persisting to a bucket solves the
durability problem that bites container platforms with ephemeral filesystems
(\S\ref{chap:build}): the data outlives any individual instance, and multiple
instances can point at the same location.

\section{URI Detection}

The feature adds scheme detection so the connector can tell a bucket URI from a
filesystem path:

\begin{lstlisting}
const OBJECT_STORE_SCHEMES: &[&str] = &[
    "s3://", "s3a://", "gs://", "gcs://",
    "az://", "azure://", "abfs://", "abfss://",
];
pub fn is_object_store_uri(location: &str) -> bool {
    OBJECT_STORE_SCHEMES.iter().any(|s| location.starts_with(s))
}
\end{lstlisting}

\section{The Extended Connector}

\verb|connect| gains a \verb|storage_options| argument and two behavioural
changes versus the \verb|main| version:

\begin{lstlisting}
pub async fn connect(
    location: &std::path::Path,
    storage_options: &std::collections::HashMap<String, String>,
) -> anyhow::Result<Connection> {
    // Only create local directories for filesystem paths.
    if !is_object_store_uri(path_str) {
        if let Some(parent) = location.parent() {
            std::fs::create_dir_all(parent).ok();
        }
    }
    let mut builder = lancedb::connect(path_str);
    if !storage_options.is_empty() {
        builder = builder.storage_options(storage_options.clone());
    }
    builder.execute().await
}
\end{lstlisting}

Two things change relative to \verb|main|: object-store URIs are passed to
LanceDB verbatim (no \verb|create_dir_all|, no path joining), and
credential/endpoint options are threaded through
\verb|builder.storage_options(...)|. Everything above the connector --- schemas,
\verb|init_db|, the migration, every route --- is byte-for-byte identical. The
encryption model and data layout do not change; only where the bytes land does.

\section{Configuration}

On the feature branch, \verb|config.yaml| grows an optional free-form
\verb|storage| section whose key/value pairs become the LanceDB storage options,
and the database path becomes a URI:

\begin{lstlisting}
database:
  path: "s3://my-bucket/chat.lance"      # or gs://, az://, abfs://

storage:
  aws_endpoint: "http://localhost:4566"  # LocalStack / MinIO / moto
  region: "us-east-1"
  allow_http: "true"
  aws_access_key_id: "test"
  aws_secret_access_key: "test"
\end{lstlisting}

Standard AWS environment variables (\verb|AWS_ACCESS_KEY_ID|,
\verb|AWS_SECRET_ACCESS_KEY|, \verb|AWS_REGION|, \verb|AWS_ENDPOINT|,
\verb|AWS_ALLOW_HTTP|) are also honoured, so no secrets need be written to disk in
production.

\section{The Smoke Test}

The branch adds \verb|tests/object_storage.rs|, a single
\verb|#[tokio::test] #[ignore]| case --- ignored by default because normal CI has
no object store. It connects to an S3-compatible endpoint, runs \verb|init_db|,
verifies all five tables are created in the bucket, and writes then reads a row
to prove a full round-trip. It reads its target from environment variables:

\begin{center}
\begin{tabular}{@{}ll@{}}
\toprule
\textbf{Variable} & \textbf{Example} \\
\midrule
\verb|SEAL_TEST_S3_URI| & \verb|s3://seal/it.lance| \\
\verb|SEAL_TEST_S3_ENDPOINT| & \verb|http://localhost:5001| \\
\verb|SEAL_TEST_S3_REGION| & \verb|us-east-1| \\
\verb|AWS_ACCESS_KEY_ID| & \verb|test| \\
\verb|AWS_SECRET_ACCESS_KEY| & \verb|test| \\
\bottomrule
\end{tabular}
\end{center}

\subsection*{Running it against \texttt{moto} (Docker-free)}

\begin{lstlisting}
# A pure-Python S3 mock -- no Docker required.
uv venv /tmp/motoenv
uv pip install --python /tmp/motoenv/bin/python "moto[server]"
/tmp/motoenv/bin/moto_server -p 5001 &

AWS_ACCESS_KEY_ID=test AWS_SECRET_ACCESS_KEY=test AWS_DEFAULT_REGION=us-east-1 \
  aws --endpoint-url http://localhost:5001 s3 mb s3://seal

SEAL_TEST_S3_URI=s3://seal/it.lance \
SEAL_TEST_S3_ENDPOINT=http://localhost:5001 \
  cargo test --test object_storage -- --ignored --nocapture
\end{lstlisting}

The same test works against LocalStack, MinIO, or real S3 by pointing the
endpoint and credentials elsewhere.

\section{A Note on Cargo Features}

Talking to cloud stores requires LanceDB's object-store backends to be compiled
in. When merging this feature, confirm the \verb|lancedb| dependency in
\verb|Cargo.toml| enables the needed backends (e.g. \verb|aws|, \verb|gcs|,
\verb|azure|); on \verb|main| the dependency is declared without those features
because the connector never reaches for them.
